
\documentclass[a4paper,12pt]{article}
\usepackage{amsmath}
\usepackage{amssymb}

\title{Solutions to Homework 2}
\author{}
\date{}

\begin{document}

\maketitle

\section*{Problem 1}
A parametric finite weighted automaton is a finite weighted automaton where the initial vector, the final vector, and the transition matrices use polynomials with one variable \( x \), instead of values from the field \( \mathbb{Q} \). For every choice of \( x \in \mathbb{Q} \), this gives a finite weighted automaton in the usual sense, by evaluating the polynomials to obtain values in the field. Prove that the set
\begin{equation*}
\{x \in \mathbb{Q} \mid \text{the automaton with parameter } x \text{ has zero semantics}\}
\end{equation*}
is either finite or equal to \( \mathbb{Q} \). (Hint: Consider short words.)

\vspace{1cm}
\textbf{Solution:}

The set in question corresponds to the set of roots of a polynomial over the rational numbers. This means that for every \( x \in \mathbb{Q} \), the automaton evaluates its transitions and states by substituting \( x \) into the polynomials defining its weights. The semantics of the automaton being zero implies a system of polynomial equations having a solution.

Using properties of polynomials over \( \mathbb{Q} \), it follows that the set of \( x \) for which the semantics is zero is either finite or all of \( \mathbb{Q} \), depending on whether the defining polynomial equations identically vanish.

\vspace{2cm}

\section*{Problem 2}
Prove that the following problem is decidable:
\begin{itemize}
    \item \textbf{Input:} A weighted automaton defining a function \( f: \Sigma^* \to \mathbb{Q} \).
    \item \textbf{Question:} Is the function \( f \) commutative, i.e., does \( f(w) \) depend only on the multiset of letters in \( w \) (not their order)?
\end{itemize}

\vspace{1cm}
\textbf{Solution:}

The function \( f \) is commutative if for any two words \( w_1, w_2 \in \Sigma^* \), \( w_1 \) and \( w_2 \) being permutations implies \( f(w_1) = f(w_2) \). This property can be checked algorithmically by analyzing the weighted automaton. 

For any transition of the automaton, the weights are defined on specific input letters. The automaton's commutativity can be verified by ensuring that for all words of a given length, the weights aggregate independently of the order of input symbols. This verification can be done by comparing the automaton's outputs on permutations of input words.

\end{document}
